%%%%%%%%%%%%%%%%%%%%%%%%%%%%%%%%%%%%%%%%%%%%%%%%%%%%%%%%%%%%%%%%%%%%%%
%%                     Catalysis
%%%%%%%%%%%%%%%%%%%%%%%%%%%%%%%%%%%%%%%%%%%%%%%%%%%%%%%%%%%%%%%%%%%%%%
%\color{blue}
\subsection{Glyph: \glyph{Catalysis}}\label{sec:catalysis}

A catalysis is a particular case of stimulation, where the effector affects
positively the flux of a process represented by the target transition. The positive effect on the transition is due to the lowering of the activation energy of a reaction.

\begin{glyphDescription}
 \item[SBO]\mbox{}\\ SBO:0000172 ! catalysis
 \item[origin]\mbox{}\\ Any EPN (section~\ref{sec:EPNs}) or any logical operator (section~\ref{sec:logic}).
 \item[target]\mbox{}\\ Any transition node (section~\ref{sec:PNs}).
 \item[node]\mbox{}\\ The target extremity of a \glyph{catalysis} carries an empty circle.
 \end{glyphDescription}

\begin{figure}[H]
  \centering
  \includegraphics[scale = 0.5]{images/catalysis}
  \caption{The \PD glyph for \glyph{catalysis}.}
  \label{fig:catalysis}
\end{figure}


